\section{A standalone asymmetric parity interface}
\label{sec:parity}
This component bypasses the inflated central-binomial index, but does
\emph{not} yet constitute a new complete consistency equation. Its
precise theorem concerns positive input integers $S,T$:
\[
 S\mathbin{\&}T=0\quad\Longleftrightarrow\quad
 \exists\mathbf w\in\N^{12}\ P(S,T,\mathbf w)=0.
\]
The bit operation occurs only in this specification, not in the polynomial.
The three-Pell quotient technique has a precedent in Jones--Sato--Wada--Wiens
(1976), Section3.9. The specialization here uses odd-quotient proximity,
rank locking, and a built-in congruence to remove one witness.

Write $n=S+T$ and use the twelve nonnegative witnesses
$y,r,c,p,l,i,j,u_0,u_1,u_2,u_3,v$. Define
\begin{align*}
 x&=2(y+1),&Y&=2r+1,\\
 M&=128(n+1)(x+1)(Y+1),&A&=M(x+1),\\
 C&=n+1+c,&K&=T+1+p(M-1),\\
 L&=S+1+l(Mx-1),&D&=(A^2-1)C^2+1,\\
 E&=2(i+1)DC^2,&F&=(A^2-1)E^2+1,\\
 G&=A+F(F-A),&H&=C+F\bigl(n+1+(2j+1)C\bigr),\\
 I&=(G^2-1)H^2+1.&&
\end{align*}
The five residuals are
\begin{align*}
 U_0&=u_0^2-\bigl((2n+6)^3(2n+8)(y+1)^2+1\bigr),\\
 U_1&=u_1^2-\bigl((M^2-1)K^2+1\bigr),\\
 U_2&=u_2^2-\bigl(((Mx)^2-1)L^2+1\bigr),\\
 U_3&=u_3^2-DFI,\\
 U_4&=K^2L^2-64(C-YKL)^2-v-1.
\end{align*}
Set $P=\sum_{a=0}^4U_a^2$. This is a literal ordinary polynomial of
exact total degree 504, counting the two input coordinates as variables.
The independent audit corrects an earlier draft's bound of 500.
A thirteen-witness, six-equation version retains a separate divisibility
witness and has degree368. These are distinct tradeoffs.

\paragraph{Pell conventions and hypotheses.}
Let $\psi_a(k)$ be the second coordinate of
$(a+\sqrt{a^2-1})^k$. For $a>1$, every nonnegative solution of
$(a^2-1)z^2+1=\square$ is such a coordinate. Also
$\psi_a(k)\equiv k\pmod{a-1}$,
$(2a-1)^{k-1}\le\psi_a(k)\le(2a)^{k-1}$, and
\[
 (2a)^q\left(1-\frac{q}{2a}\right)
 \le\psi_a(q+1)\le(2a)^q.
\]
The growth-square theorem applied to $U_0$ gives $x>128\cdot2^n$.
The indexed-Pell certificate of JSWW Lemma3.8 gives
$C=\psi_A(n+1)$ from $C\ge n+1>1$, $DFI=\square$, and the
prescribed two congruences of $H$. Here $F\equiv1\pmod{2C}$, so
our formula for $H$ supplies $H\equiv n+1\pmod{2C}$ and
$H\equiv C\pmod F$, with $H\ge n+1+2C$.

\paragraph{Soundness and rank locking.}
The Pell equations and congruences give
\[
 K=\psi_M\bigl(T+1+a(M-1)\bigr),\quad
 L=\psi_{Mx}\bigl(S+1+b(Mx-1)\bigr),\qquad a,b\ge0.
\]
The last residual implies $|C/(KL)-Y|<1/8$ and therefore
$C/(KL)>7/8$. If $a>0$, the lower growth bound gives
$K\ge(2M-1)^{2n+1}>2(2A)^n\ge2C$, because
$M-1\ge2n+1$ and $(2M-1)^2>2A$. The same argument applies to
$L$ if $b>0$. Either case contradicts the proximity inequality.
Thus both ranks are the intended ones.

Set $\rho=(x+1)^n/x^S$, $\sigma=C/(KL)$, and
$\epsilon=n/(2M)$. The relative Pell bounds place both normalized
numerator and denominator in $[1-\epsilon,1]$, whence
$|\rho-\sigma|\le\sigma\epsilon/(1-\epsilon)<1/128$.
The binomial expansion gives $\rho=Q+f$ with integer
$Q\equiv\binom nS\pmod x$ and $0\le f<2^n/x<1/128$.
Consequently $|Y-Q|<1/8+1/128+1/128<1$. Hence $Y=Q$ is odd.
Since $x$ is even, the binomial coefficient is odd. Lucas' binary
criterion yields $S\mathbin{\&}T=0$.

\paragraph{Completeness and the congruence compression.}
For disjoint positive $S,T$, choose a growth-square witness $y$ and
put $Y=\lfloor(x+1)^n/x^S\rfloor$, an odd positive integer. Define
$M,A,C,K,L$ at their intended Pell ranks. Their congruences supply
nonnegative $c,p,l$. The indexed-Pell theorem supplies $i$ and an
initial coordinate $H_{\rm old}$ satisfying both congruences.
Its coprimality identities separate $D,F,I_{\rm old}$ into squares,
so $H_{\rm old}$ is a Pell coordinate for base $G$.
The determinant-one Pell matrix is periodic modulo $2CF$.
Increasing its index by multiples of a positive period preserves both
congruences and gives arbitrarily large $H_{\rm new}$. Choose it at
least $C+F(n+1+C)$. Writing $H_{\rm new}-C=Fq$ yields
$q\equiv n+1+C\pmod{2C}$, hence
$q=n+1+(2j+1)C$ with $j\ge0$. This proves completeness of the
compressed formula for $H$, rather than assuming an unproved bound
for a particular canonical witness.
The same relative estimates give $|\sigma-Y|<1/64<1/8$, so the
last positive integer slack exists. The square-root witnesses exist
by construction. This completes both directions of the theorem.

\paragraph{Exact degree and boundary.}
The successive degrees of $M,A,C,D,E,F,G,H,I$ are
$3,4,1,10,13,34,68,36,208$. Thus $DFI$ has degree252 and
$U_3^2$ degree504, with a nonzero leading term; no other squared
residual cancels it. The theorem assumes $S,T>0$ and all twelve
witnesses nonnegative. Substituting a packed source would still need
a separate proved recovery of the binary base. That missing step is
not hidden in this theorem, and no completed global improvement is
claimed for this component.
